\scanpage{150}
\unsure{contragrédience}, sur $V(M)$].

Une structure de coalg $\Delta$ sur $M$ est dite \uline{compatible} avec la graduation, si
\[
\Delta:M\mathbin{\text{---}}M\otimes M
\]
est compatible avec la ``graduation totale''
\[
\left(\Delta(M_n)\subset\sum_{i+j=n}M_i\otimes M_j\right)
\]
ou, ce qui revient au $=$, si la loi de composition sur $V(M)$ est compatible avec les graduations, i.e. provient d'hom
\diagram{
\[
\begin{array}{ccccc}
V(M_i)&\times&V(M_j)&\longrightarrow&V(M_{i+j})\\
\text{\textquotedbl}&&\text{\textquotedbl}&&\text{\textquotedbl}\\
V(M)^i&\times&V(M)^j&&V(M)^{i+j}
\end{array}
\]
}
Il revient aussi de $=$ de dire que la multiplication dans $V(M)$ (ou la comultiplication de $M$) commute à l'action de $G$.
\margin{La donnée d'une telle $\Delta$ revient à la donnée des
\[
\Delta_n:M_n\mathbin{\text{---}}\sum_{i+j=n}M_i\otimes M_j
\]
i.e. de \deleted{$\unsure{\sum}$}
\diagram{\begin{tikzcd}
\displaystyle\prod_{i+j=n}\underbrace{V(M_i)\times V(M_j)}_{V(M)^i\times V(M)^j}\arrow[r] & V(M_{i+j})\arrow[d,phantom,"\Vert" description]\\
& V(M)^n
\end{tikzcd}}
}
\margin{\uline{NB} dans ce pt de vue ``familles des modules indexés par $\unsure{M_i}$'', au lieu des $\Delta_n$ il suffit de se donner des
\[
\Delta^{ij}:M_{i+j}\mathbin{\text{---}}M_i\otimes M_j,
\]
ou des accouplements
\[
V(M)^i\times V(M)^j\mathbin{\text{---}}V(M)^{i+j},
\]
\unsure{sans} condition loc. finitude. \unsure{on a des algèbres}\illegible{2}
}
Bien sûr, les (morph. de coalgèbres graduées) correspondent aux hom des $O_k$-Alg. correspondantes (éventuellement graduées).

Condition d'anticommutativité, lorsqu'on dispose d'un hom
\[
\deleted{\unsure{\varphi}}\quad Z\mathbin{\text{---}}\mathbb{Z}/2\mathbb{Z}
\]
(i.e. $\mu_2\mathbin{\text{---}}G$) qui permet de \deleted{\unsure{d'introduire à $M$}}\inserted{d'associer à la graduation donnée} une graduation de \deleted{de} groupe $\mathbb{Z}/2\mathbb{Z}$, idem pour $V(M)$,
\margin{149}
